\section{Method and implementation details}
\label{app:forge}

This appendix specifies program and graph encoding, model training, sampling, verification and
route assessment. Algorithm~\ref{alg:forge-inference} combines the inference steps.

\subsection{Reaction-program encoding}
\label{app:program-encoding}

The program tuple in Equation~\ref{eq:reaction-program} uses the ordered atom-position set
$\mathcal I_{\rm atom}(P)=\{1,\ldots,n(P)\}$. The arrays $\mathbf r$ and $\mathbf k$ assign
precursor roles and core positions to these atoms. This atom-position index differs from
$\mathcal I(P)$, which indexes all six families of categorical coordinates.

Morphology is computed from the selected atom-mapped trace and its sparse serialization.
For role $r$, let $V_r^{\rm ext}$ contain its exterior atoms, excluding reaction-core atoms, and
let $o_r(v)$ count the same-role exterior children of $v$ in the serialized tree. The four fields are
\begin{equation}
 M_r=(n_r,j_r,\gamma_r,a_r),\qquad
 n_r=|V_r^{\rm ext}|,\qquad
 j_r=\sum_{v\in V_r^{\rm ext}}\max\{o_r(v)-1,0\}.
\end{equation}
Here $\gamma_r$ counts closure edges with both endpoints in $V_r^{\rm ext}$, and $a_r$ counts
same-role bonds between exterior and core atoms. When the serialized tree restricts to a spanning
forest of the exterior subgraph, $\gamma_r$ equals that subgraph's cycle rank. The junction budget
is a tree-branching count; it does not count degree contributed by closure bonds. Consequently these
fields depend on the selected annotation and serialization, not on the unannotated graph alone.
Fixed-coordinate equality alone does not establish agreement with every morphology field.
Positions without precursor origin use the designated absent morphology state. Integer morphology
values are encoded with an offset of one, reserving zero for that absent state.

The training-fold program prior is defined in Section~\ref{sec:program-conditioning}.
The memory encoder uses separate program-identity and depth tokens,
$m_c=e_c(c)+e_{\rm type}(0)$ and $m_d=e_d(d)+e_{\rm type}(1)$, together with one token per active
atom position,
\begin{equation}
 m_i=e_r(r_i)+e_k(k_i)+e_{\rm pos}(i)+e_{\rm rep}(g_i)+e_{\rm comp}(u_i)
     +\sum_{v=1}^{4}e_{M,v}(M_{r_i,v})+e_{\rm type}(2).
\end{equation}
Here $g_i$ is the repeat-group state and $u_i$ the position within the precursor-origin component;
they are structural conditioning variables, not precursor identity labels. All embeddings share the
model hidden dimension. Layer normalization precedes memory self-attention, followed by residual
attention and feed-forward updates. The resulting program, depth and atom-position tokens provide
the memory for cross-attention in the graph denoiser. The morphology fields determine active
positions and fixed or variable masks. Sampling from the prior specifies these constraints;
generation assigns the atom, bond and topology states within each role.

A chemistry-specific adapter determines the precursor roles, reaction core, reverse decomposition
and forward transform associated with $c$. Its sampling layout determines active positions and
fixed or variable coordinates; the embeddings condition neural prediction on the remaining
program fields.

\subsection{Graph serialization and source annotations}
\label{app:graph-encoding}

A spanning-tree representation with at most twelve residual closure edges preserves the full
254-heavy-atom support without an $O(N_{\max}^2)$ edge state. The discrete flow generates the atom, bond
and topology states in this representation.
Let $\alpha=(\mathbf C,\tau,\mathbf r,\mathbf k)$ contain the selected exact precursor tuple,
registered assembly trace, atom-wise precursor-role assignment and reaction-core annotation. The deterministic
program-aware encoder
\begin{equation}
    \operatorname{Enc}_P:
    \{(G,\alpha):G\in\mathfrak G_{N_{\max},K_{\rm cl}}
      \text{ and }\alpha\text{ is admitted for }P\}
    \longrightarrow \mathfrak X_{N_{\max},K_{\rm cl}}
\end{equation}
constructs the categorical graph state
\begin{equation}
    X_\star:=\operatorname{Enc}_P(G,\alpha)=
    \left(
        \mathbf{a},
        \boldsymbol{\pi},
        \mathbf{b},
        \mathbf{e}^{L},
        \mathbf{e}^{R},
        \mathbf{z}
    \right)
\end{equation}
where $\mathbf{a}$ contains atom states, $\boldsymbol{\pi}$
parent pointers, $\mathbf{b}$ spanning-tree bond states, $\mathbf{e}^{L}$ and $\mathbf{e}^{R}$
closure-endpoint pointers, and $\mathbf{z}$ closure-bond states.

The encoder first partitions the atoms into connected precursor-origin components under $\alpha$.
It selects as root the reaction-core atom with the smallest tie-broken canonical RDKit rank, and
orders the origin-component graph by breadth-first traversal from that root component, breaking
ties by role index, component canonical ranks and original atom index. Within each component it uses
a deterministic depth-first preorder from the selected attachment atom, with neighbors ordered by
the same canonical ranks. The traversed bonds define the spanning tree. Every residual bond is
represented once as an ordered endpoint pair with its bond label, and residual triples are sorted
lexicographically. Role, core, fixed-coordinate and morphology arrays follow the
same atom order. Each admitted source record supplies a selected qualified exact assembly trace
and its precursor-origin assignment. Alternative source-role bindings remain explicit annotations;
they do not receive additional graph sampling mass when their product constitutions coincide. Thus
$\operatorname{Dec}(\operatorname{Enc}_P(G,\alpha))\cong G$ for every admitted record.

The raw padded categorical product space $\overline{\mathfrak X}_{N_{\max},K_{\rm cl}}$ also
contains malformed states. Its well-formed subset $\mathfrak X_{N_{\max},K_{\rm cl}}$ enforces an
active node prefix, earlier-parent pointers, distinct valid closure endpoints, no duplicate edges,
and null values in unused closure slots. The total decoder, including its failure value, is typed as
\begin{equation}
    \operatorname{Dec}:\overline{\mathfrak X}_{N_{\max},K_{\rm cl}}
    \longrightarrow
    \mathfrak G_{N_{\max},K_{\rm cl}}\cup\{\bot_{\rm dec}\}.
\end{equation}
For this study, $N_{\max}=254$ and $K_{\rm cl}=12$. The atom vocabulary includes qualified bromine states; bond states encode single, double and triple bonds, with aromatic structures represented through Kekul\'e bond assignments.
These bounds specify architectural support. Reaction consistency, nonzero probability under a
trained model and synthesis-program completeness require separate conditions.

The structural decoder specifies graph reconstruction from an already chosen state.
The production terminal readout additionally chooses that state from neural predictions
(Appendix~\ref{app:sampling}). These are separate operations.

The chemistry adapter marks any atom, pointer, or bond states fixed by $P$.
FORGE holds these states constant while corrupting the remaining states along the linear probability
path and predicting their clean values from $(X_t,t,P)$. Let $\Omega_P$ be the subset of the raw state
space satisfying the allowed coordinate alphabets and fixed values, and define the admissible program set
\begin{equation}
    \mathcal P_{\rm adm}
    =\left\{P:\mathfrak X_{N_{\max},K_{\rm cl}}\cap\Omega_P\neq\varnothing\right\}.
\end{equation}
The morphology prior is normalized only on $\mathcal P_{\rm adm}$. Each chemistry adapter is required
to satisfy the program-compatible encoding property: every admitted annotated record obeys
$\operatorname{Enc}_P(G,\alpha)\in\Omega_P$. Corpus admission requires this property, an exact
constitutional encode--decode round trip and source-consistent origin/core assignments. Records
failing admission are excluded. The frozen split is \PendingValue{split manifest}; noise marginals,
program priors and motif references are fitted only to its training partition. These checks do not
establish the property for unregistered chemistry.

\subsection{Transformer architecture and training objectives}
\label{app:training}

The denoiser is a graph Transformer over the complete sparse state. Each layer applies graph
self-attention with relation biases for parent, child and closure edges, then cross-attends to memory
tokens representing program identity, program depth, precursor role and reaction-core position.
Cross-attention is repeated in every layer so that reaction semantics can influence both local state
updates and long-range coordination among precursor-derived regions. The backbone has six layers,
hidden dimension 192, eight attention heads and dropout 0.1. Each routed adapter mixes three
experts with bottleneck dimension 64 using weights computed from the program summary. Additional
routed adapters feed atom/parent-bond predictions and closure predictions. Auxiliary heads predict
precursor roles, core positions and exterior child counts. The shared backbone receives
repeat-group, component-position and role-morphology conditioning.

Training assigns equal mass to the 22 admitted formal families and uniform mass to qualified
unique product constitutions within each family. Sampling is with replacement; multiple source
bindings do not multiply a constitution's probability, and sampling weights are not applied again
to the loss. Each optimizer update contains 3,168 examples from three distinct families, with
1,056 examples per selected family. A seeded permutation defines a balanced 22-update cycle in
which every family appears three times. Family batches are sharded across eight H100 GPUs, and
their coordinate reductions are accumulated before family-gradient projection. Padding is trimmed
to the longest molecule within a family batch while retaining the full closure support.

The optimizer is AdamW with learning rate $10^{-4}$, weight decay $5\times10^{-5}$,
$(\beta_1,\beta_2)=(0.9,0.999)$, $\epsilon=10^{-8}$ and global gradient-norm clipping at 1.
The update budget is \PendingValue{optimizer updates}; the independent training seeds and
checkpoint-selection rule are \PendingValue{training seeds and checkpoint rule}. Model and gradient
computations use deterministic FP32 arithmetic. Checkpoints retain model, optimizer, random-number
streams and the committed sampler state; uncommitted prefetch work is discarded on restart.
Training times are sampled uniformly, $t\sim\mathcal U(0,1)$, without endpoint clipping.

In the task losses below, $c$ indexes a formal family; multiple registered program variants in that
family share its task reduction.
For formal family $c$ and state family
$q\in\mathscr J=\{\mathbf a,\boldsymbol\pi,\mathbf b,\mathbf e^L,\mathbf e^R,\mathbf z\}$,
let $\mathcal S_{c,q}$ be the selected variable coordinates across that family's minibatch.
With coordinate loss $\ell_i=-\log d_{\theta,i}(X_{\star,i}\mid X_t,t,P)$, the pooled reduction is
\begin{equation}
 L_{c,q}^{\rm pool}=\frac{\sum_{i\in\mathcal S_{c,q}}\ell_i}{\max\{1,|\mathcal S_{c,q}|\}}.
\end{equation}
Coordinates here include their example index. Chemistry states (atoms, parent bonds and closure
bonds) additionally use equal mass over roles present in the selected coordinates. If
$\mathcal S_{c,q,r}$ contains coordinates assigned to role $r$ and
$\mathcal R_{c,q}=\{r:|\mathcal S_{c,q,r}|>0\}$, this reduction is
\begin{equation}
 L_{c,q}^{\rm role}=\frac{1}{|\mathcal R_{c,q}|}
 \sum_{r\in\mathcal R_{c,q}}\frac{\sum_{i\in\mathcal S_{c,q,r}}\ell_i}{|\mathcal S_{c,q,r}|},
 \label{eq:role-balanced-loss}
\end{equation}
with zero loss if no role is present. Atom and parent-bond groups use their position's role;
closure-bond groups use the role at the clean left endpoint. Pointer losses use the pooled reduction
and mask logits to the permitted layout positions. The flow term $L_{c,\rm flow}$ sums these six
reduced losses, with unit weights.

For each family, the task loss is
\begin{equation}
\begin{aligned}
 L_c={}&L_{c,\rm flow}
 +0.25L_{c,\rm role}+0.25L_{c,\rm core}+0.25L_{c,\rm repeat}\\
 &+L_{c,\rm child}+0.5L_{c,\rm junction}+L_{c,\rm chem\mid top}.
\end{aligned}
\label{eq:production-task-loss}
\end{equation}
The role and core terms are cross-entropies averaged equally over present target states.
The repeat term averages squared differences between predicted atom distributions, and separately
parent-bond distributions, over exact retained-fragment correspondences in repeated precursor
instances; it sums the two means. Atom correspondence requires a unique labelled fragment match
that preserves core annotations. Parent-bond correspondence matches both edge endpoints rather
than equating serialized child positions. Unequal fragments or ambiguous matches contribute no
repeat pairs; these target-derived alignments supervise the loss and are not supplied at inference.
The child term is state-balanced cross-entropy for same-component exterior child
counts in $\{0,1,2,3\}$. The junction term applies smooth-L1 loss to predicted and target
role-local junction totals, each divided by the number of exterior atoms in that role, then averages
over eligible examples within a role and equally across present roles. Predicted junction totals
sum the expected child-count contribution $\max\{o-1,0\}$ across exterior atoms.
The topology-conditioned chemistry term repeats the role-balanced atom and bond losses in a second
forward pass with clean parent and closure endpoints and corrupted atom and bond states.
Each ablation changes only the declared objective or architecture terms under its matched training contract.

PCGrad combines the task gradients in integer formal-family-state order. For
$g_c=\nabla_\theta L_c$, it initializes $\widetilde g_c=g_c$ and, for each other task $h$ in that
order, applies
\begin{equation}
 \widetilde g_c\leftarrow\widetilde g_c-
 \mathbf1\{\langle\widetilde g_c,g_h\rangle<0\}
 \frac{\langle\widetilde g_c,g_h\rangle}{\max\{\|g_h\|_2^2,10^{-12}\}}\,g_h
\end{equation}
on parameter coordinates available to both tasks. Gradients are not norm-normalized, and the
optimizer receives the arithmetic mean of the projected task gradients. This update is not generally
the gradient of a scalar sum of task losses. No population optimality claim is made for the full
procedure. In particular, target-dependent group reductions and the auxiliary objectives require
analysis beyond Proposition~\ref{prop:masked-bayes}.

Layerwise cross-attention, auxiliary losses, routed adapters and gradient-conflict control are
separate interventions in the architecture comparison (Table~\ref{tab:22-architecture}).

\subsection{Numerical sampling and terminal decoding}
\label{app:sampling}

The sampler uses \PendingValue{flow steps} steps, denoted $L_{\rm samp}$. Its strictly positive atom and bond
source marginals are indexed by registered program and precursor role and fitted to TRAIN data.
Pointer marginals are restricted and normalized on their allowed layout positions. The finite
sampler evaluates the denoiser at $t_k=k/L_{\rm samp}$, for
$k=0,\ldots,L_{\rm samp}-1$, and uses $\Delta t=1/L_{\rm samp}$.
Conditional on $Z_k=x$, it draws one clean category
$A_{k,i}\sim d_{\theta,i}(\cdot\mid x,t_k,P)$ independently for each variable coordinate.
Write $s_{t,i}^P(u\mid a)=(1-t)q_{0,i}^P(u)+t\mathbf1\{u=a\}$ for its conditional path factor,
$\dot s_i^P(u\mid a)=\mathbf1\{u=a\}-q_{0,i}^P(u)$ for the derivative, and
$[z]_+=\max\{z,0\}$. The code uses the conditional off-diagonal rate
\begin{equation}
 R_{t,i}^{*,\varepsilon}(u,v\mid a,P)
 =\frac{[\dot s_i^P(v\mid a)-\dot s_i^P(u\mid a)]_+}
 {|\mathcal X_i(P)|\max\{s_{t,i}^P(u\mid a),\varepsilon\}},
 \quad u\neq v,\qquad \varepsilon=10^{-8}.
 \label{eq:clamped-rate}
\end{equation}
Full source support is required only on the allowed coordinate alphabet, including the masked
candidate set for a pointer. Thus that alphabet is the support for every $t_k<1$.
The denominator clamp in Equation~\ref{eq:clamped-rate} is a numerical safeguard absent from the
ideal rate in Appendix~\ref{app:flow-proof}.

For the sampled endpoint $a=A_{k,i}$, let
$h_{k,i}=\Delta t\sum_{v\neq x_i}R_{t_k,i}^{*,\varepsilon}(x_i,v\mid a,P)$ and set
$\kappa_{k,i}=\min\{1,0.999/h_{k,i}\}$ when $h_{k,i}>0$, with $\kappa_{k,i}=1$ otherwise.
The implemented coordinate kernel is
\begin{equation}
 \widehat K_{k,i}(v\mid x,a)=
 \begin{cases}
 \kappa_{k,i}\Delta t\,R_{t_k,i}^{*,\varepsilon}(x_i,v\mid a,P),&v\neq x_i,\\
 1-\kappa_{k,i}h_{k,i},&v=x_i.
 \end{cases}
 \label{eq:capped-kernel}
\end{equation}
Its off-diagonal masses sum to at most $0.999$; the diagonal is the remaining mass. Consequently
all entries are nonnegative and each row sums to one. Coordinates are drawn independently from
these kernels, followed by the projection $\Pi_P$.

Without the clamp or cap, and when every conditional Euler row is nonnegative, averaging over
$A_{k,i}$ gives the Euler kernel built from Equation~\ref{eq:posterior-rate} exactly. One sampled
endpoint per coordinate therefore does not by itself bias that uncapped one-step law. The nonlinear
cap generally changes the averaged kernel. Independent coordinate updates also permit simultaneous
finite-step changes, so their product is a numerical approximation to the local CTMC, not its exact
transition kernel. No discretization error bound is asserted near the terminal singularity.

After all steps, the model evaluates $(Z_{L_{\rm samp}},1,P)$ once more. The production terminal
readout $H_P^\theta$ selects coordinate values using constrained argmax, restores fixed states and
returns either a raw state or a typed terminal failure, denoted collectively by $\bot_{\rm term}$.
The atom-aware readout first chooses atom states subject to unavoidable incident-bond demands
from the fixed core and non-root tree attachments; subsequent variable bonds must fit those atom
capacities. It abstains when this bounded assignment cannot be completed. It does not exchange an
element for a higher-valence element merely to accommodate a chosen topology. The registered
core condition reserves the required bond and hydrogen states for each program. In the AGILE-type
Ugi 3-CR, these include one hydrogen on the aldehyde-derived $\alpha$ carbon and one on the
isocyanide-derived amide nitrogen; these conditions are not imposed on other Ugi variants.
The selected state is then reconstructed and checked for
connectedness, valence and sanitization. A readout failure retains its detailed reason in the ledger;
for graph-law notation we extend $\operatorname{Dec}(\bot_{\rm term})=\bot_{\rm dec}$.
If $q_{\theta,\rm out}^P$ is the raw law after terminal readout, including failure, then
\begin{equation}
 p_\theta(\cdot\mid P)=\operatorname{Dec}_{\#}q_{\theta,\rm out}^P,
 \qquad
 p_\theta(G\mid P)=\sum_{z:\operatorname{Dec}(z)=G}q_{\theta,\rm out}^P(z).
 \label{eq:generated-law}
\end{equation}
This law retains structural failures and chemically invalid graphs before the validity and L1
filters; it is not conditioned on acceptance.
The ideal theorem does not include the learned $t=1$ readout, which is unconstrained by a loss
integrated over $t\in(0,1)$. Decoder and source-prior comparisons use their declared fixed-weight
or retrained contrasts (Table~\ref{tab:22-decoder}).

\subsubsection{Checked construction and selection}

Let $\Pi$ specify a bounded inference policy: the draw count, ordered proposal operators, TRAIN
references, search limits, source verifier, selection objective and selection-group size. Its
version is \PendingValue{inference policy}; the draw count $D$ is \PendingValue{draws per request},
the stage caps are \PendingValue{proposal and search caps}, and the selection-group size is
\PendingValue{selection batch size}. For request $b$, retain all raw trajectories and form
\begin{equation}
 \mathcal C_b=\{G^{\rm raw}_{b,d}:1\leq d\leq D\}
 \cup\mathcal R_\Pi\!\left(P_b,\{Z^{\rm out}_{b,d},\lambda_{b,d}\}_{d=1}^{D}\right),
\end{equation}
where $\lambda_{b,d}$ denotes saved terminal predictions and $\mathcal R_\Pi$ is the bounded
construction map. Original graphs, failed trajectories, rejected proposals and exceptions remain
in the attempt ledger. A request is not removed when its candidate set has no eligible graph.

Construction can tie repeated components, restore source-required component constraints, propose
qualified retained-head sites, and assign TRAIN-supported joint atom/bond or ring-system motifs.
The operators act on generated atom slots while preserving the declared atom budget, precursor
origins and fixed reaction core. Atom relabelling and topology edits are recorded separately.
Ring repair is applied after component and head construction, and all checks are evaluated again
on the final graph. Repeated arms receive tied edits when their established correspondences permit
this; every proposed product must still pass the unchanged source verifier. A TRAIN motif is an
explicit inference prior and does not itself certify chemical plausibility or a complete route.
An unsupported or context-dependent ring fragment is unassessed rather than assigned a fabricated
known identity. Graph cycle rank, ring systems, SSSR rings and serialization-dependent fundamental
cycles remain distinct quantities.

The selector maps the request-indexed pools to at most one graph per request,
\begin{equation}
 (G^{\rm sel}_1,\ldots,G^{\rm sel}_B)
 =\mathcal S_\Pi\bigl((P_b,\mathcal C_b)_{b=1}^{B}\bigr),
 \qquad G^{\rm sel}_b\in\mathcal C_b\cup\{\bot_{\rm dec}\}.
\end{equation}
Source-exact eligibility, applicable design predicates and product/role-component diversity and
novelty constraints are explicit parts of $\Pi$. A combined design pass requires every applicable
qualified predicate to be assessed and pass; unassessed checks cannot establish that endpoint.
The selection objective and constraint reference are \PendingValue{selector objective and floors}.
Each output retains its request, draw and proposal lineage. Raw and selected yields both use the
original $B$ requests as denominator, while cost accounting includes $BD$ trajectories, all
construction attempts, source checks and selection work.

The selected joint law is the pushforward of the raw trajectory law through
$\mathcal R_\Pi$ and $\mathcal S_\Pi$. It is distinct from Equation~\ref{eq:generated-law} and
can couple different requests. The ideal flow theorem does not certify these inference maps;
independent-draw accepted-count formulas apply only when their independence assumptions hold.
A graph-changing repair changes the proposal distribution, whereas rejecting completed graphs
alone does not.

\subsection{Exact assembly verification}
\label{app:verification}

For the chemistry $c$ specified by $P$, FORGE decomposes each valid graph into role-assigned
precursors and applies the registered forward transform. Exact replay checks assembly existence
under $c$; it does not test full agreement with the sampled conditioning layout.

\begin{definition}[Exact assembly and recovered witnesses]
\label{def:exact-assembly}
For the registered forward operator $\mathcal F_c$ on precursor inputs $\mathfrak C_c$, define
\begin{equation}
 \mathcal E_c^\star(G)=\{\mathbf C\in\mathfrak C_c:G\in\mathcal F_c(\mathbf C)\},\qquad
 \mathrm{L1}^\star(G,c)=\mathbf1\{\mathcal E_c^\star(G)\neq\varnothing\}.
\end{equation}
The bounded reverse procedure returns candidate pairs $(\mathbf C,\tau)$ in $\mathcal D_c(G)$.
Equation~\ref{eq:main-exact-traces} retains the pairs passing exact forward replay, and
\begin{equation}
 \widehat{\mathrm{L1}}(G,c)=\mathbf1\{\widehat{\mathcal E}_c(G)\neq\varnothing\}.
\end{equation}
Both predicates describe assembly existence; molecular validity is checked separately.
\end{definition}
The set $\mathcal E_c^\star(G)$ contains precursor inputs, whereas
$\widehat{\mathcal E}_c(G)$ contains input--trace pairs. Their compatible comparison is
\begin{equation}
 \{\mathbf C:\exists\tau,\ (\mathbf C,\tau)\in\widehat{\mathcal E}_c(G)\}
 \subseteq\mathcal E_c^\star(G),\qquad
 \widehat{\mathrm{L1}}(G,c)\leq\mathrm{L1}^\star(G,c).
\end{equation}
The inclusion follows directly from exact replay. Its converse requires complete reverse search,
which has not been established. Throughout the paper, verified exact-L1 yield counts valid
attempts satisfying the recovered predicate. Candidate traces are compared after the registered
role- and step-aware canonicalization, so duplicate encodings do not create additional witnesses.

Post-hoc filtering changes which completed samples are retained but cannot change the proposal
distribution's exact-L1 acceptance probability.
Under independent sampling, this probability determines the complete distribution of accepted
counts under a fixed attempt budget and the expected cost of obtaining a fixed number of accepted
products (Proposition~\ref{prop:posthoc-budget}).
We distinguish decomposition coverage from forward-replay precision because a candidate disconnection
may fail to reconstruct $G$.
Candidate decomposition coverage is the fraction of valid graphs for which
$\mathcal D_c(G)\neq\varnothing$. Candidate yield additionally uses all original requests as
its denominator; undefined conditional ratios are not replaced by zero. For a declared sequence of valid evaluation graphs $G_1,\ldots,G_m$, candidate replay precision is
\begin{equation}
 \mathrm{ReplayPrecision}
 =\frac{\sum_{b=1}^m|\widehat{\mathcal E}_c(G_b)|}
 {\sum_{b=1}^m|\mathcal D_c(G_b)|},
\end{equation}
with N/E when the denominator is zero. Replay of tuples already in
$\widehat{\mathcal E}_c(G)$ has $100\%$ success by construction. This invariant does not estimate
candidate replay precision. We separately define
\begin{equation}
 \mathrm{CandidateAmbig}(G,c)=\mathbf 1\{|\mathcal D_c(G)|>1\},\qquad
 \mathrm{ExactAmbig}(G,c)=\mathbf 1\{|\widehat{\mathcal E}_c(G)|>1\}.
\end{equation}
A distinct exact-L1 product is a deduplicated generated constitution. An unambiguous exact
decomposition satisfies $|\widehat{\mathcal E}_c(G)|=1$. We distinguish product uniqueness from
decomposition uniqueness throughout.

\subsection{Recursive route assessment and dossier admission}
\label{app:route-assessment}

\subsubsection{Assessment interface}

For a valid generated graph $G$ and reaction program $P$ with chemistry $c$, route assessment
returns either a
complete synthesis dossier or a typed abstention.
Let $\mathfrak B$ denote the frozen route contract: depth and query budgets, registered transforms,
evidence-admission rules and a dated supplier snapshot. Let $\mathfrak D$ be the space of candidate
dossiers, each containing a synthesis tree and its reaction and terminal evidence, and let
$\mathfrak F_{\mathfrak B}$ be the disjoint finite set of typed failures.
Generation and assessment are defined by
\begin{equation}
    Y=\mathsf{Assess}_{\mathfrak B}(G,P,\mathbf C_T,\tau_T)
    \in\mathfrak D\cup\mathfrak F_{\mathfrak B}.
\end{equation}
$\mathsf{Assess}_{\mathfrak B}$ is invoked after the verifier returns one admitted root trace
$(\mathbf C_T,\tau_T)\in\widehat{\mathcal E}_c(G)$. It verifies the final assembly, searches for
preparation routes to unavailable precursors and checks terminal supplier evidence. $P$ and the assessor act at different stages: $P$ guides generation, whereas the assessor
constructs a dossier from the completed graph without modifying it.
FORGE accepts $(G,T)$ only when $G$ passes molecular and exact-assembly verification and $T$ traces
every required precursor to an available starting material; otherwise, it abstains.

\subsubsection{Recursive expansion and completion checks}

Completing a synthesis dossier requires evidence for every recovered precursor. Exact L1
verification specifies the precursor structures; availability and preparation routes require
additional evidence. Each precursor must therefore match an available material or have a supported
route from available starting materials.
Let $T$ be a candidate dossier with bipartite synthesis tree $(V_M,V_R,E_T)$ rooted at $G$, where each product molecule
points to its preparation reaction and each reaction points to its required reactants.
The assessor constructs these routes recursively and withholds a complete dossier when any branch
remains unresolved.
A molecule node terminates when a current, dated supplier record lists an exact constitutional match
as available; otherwise, FORGE must expand that node through a supported, forward-verified reaction.
FORGE uses a bounded level-two (L2) planner to propose each expansion.
For each proposal, FORGE records the reaction, reactants, product, and supporting evidence, and
accepts the step only when the forward transform reconstructs the target precursor exactly.
Each accepted reaction adds its reactants as new molecule nodes, and FORGE applies the same
termination-or-expansion rule recursively.
The current assessor admits only an unambiguous verified exact-L1 root. Every dossier therefore
records the sole selected root trace
$(\mathbf C_T,\tau_T)\in\widehat{\mathcal E}_c(G)$; an ambiguous exact set produces a typed
abstention. For each upstream reaction node $r$, let
$\mathrm{Admitted}_{\mathfrak B}(r)=1$ when its reaction and substrate scope satisfy the frozen
evidence policy, and let $\mathrm{Verify}_{\mathfrak B}(r)=1$ when
$\operatorname{target}(r)\in\mathcal F_r(\operatorname{reactants}(r))$.
For each leaf molecule $\ell$, let $\mathrm{Available}_{\mathfrak B}(\ell)=1$ when it satisfies the terminal
supplier-evidence rule.
Let $V_R^{\mathrm{up}}(T)$ denote the upstream preparation-reaction nodes and $L(T)$ the leaf
molecules. $\mathrm{WellFormed}(G,T,P)$ requires that $T$ is a connected, acyclic bipartite tree
rooted at $G$; each internal molecule has exactly one recorded producing reaction; no molecular constitution repeats along an ancestor path; each reaction's
molecule children equal its recorded reactants with their multiplicities; and $T$ has no disconnected, cyclic or
extraneous node. Its distinguished final-assembly record must have product $G$ and precursor
children exactly $\mathbf C_T$ under trace $\tau_T$.
\begin{equation}
\begin{aligned}
    \mathrm{Complete}_{\mathfrak B}(G,T,P)
    ={}&\mathrm{Valid}(G)\,\mathrm{WellFormed}(G,T,P)\\
    &\times\mathbf1\!\left[\widehat{\mathcal E}_c(G)=\{(\mathbf C_T,\tau_T)\}\right]\\
    &\times\prod_{r\in V_R^{\mathrm{up}}(T)}
    \mathrm{Admitted}_{\mathfrak B}(r)\mathrm{Verify}_{\mathfrak B}(r)
    \prod_{\ell\in L(T)}\mathrm{Available}_{\mathfrak B}(\ell).
\end{aligned}
\end{equation}
All factors are binary and empty products equal one. Repeated use of a material is represented by
separate tree occurrences; equality of their constitutions does not merge their nodes.
Under finite search bounds and terminating evidence queries, the assessor terminates and returns a
dossier only when $\mathrm{Complete}_{\mathfrak B}(G,T,P)=1$
(Lemma~\ref{lem:planner-termination} and Invariant~\ref{inv:dossier-soundness}).
An abstention means that FORGE did not complete a dossier under $\mathfrak B$; it does not imply that
no synthesis route exists outside the frozen search space or evidence snapshot.
The guarantee is computational and does not establish experimental synthesis success.

\subsubsection{Recorded evidence and support tiers}

For each accepted pair $(G,T)$, the synthesis dossier contains every molecule and reaction in $T$,
the evidence for each reaction and a dated supplier record for each terminal material.
We classify assembly support relative to the fixed reference space into four tiers:
\textbf{E0 (enumerated):} the generated lipid matches a product that was already constructed
computationally from known components using a supported reaction.
\textbf{E1 (known components):} the generated lipid was not included in the bounded E0 enumeration,
but a supported reaction assembles it entirely from known components.
\textbf{E2 (generated precursor):} a supported reaction assembles the lipid, but at least one
precursor is newly generated and requires a complete L2 route to available starting materials.
\textbf{E3 (new assembly family):} the lipid requires a final assembly reaction outside the
qualified reaction registry.

\subsection{End-to-end inference procedure}
\label{app:inference}

Algorithm~\ref{alg:forge-inference} records every requested output and every trajectory used to
construct it, including failures in molecular validation, exact replay or route completion.

\begingroup
\setlength{\parskip}{0pt}
\hrule height 0.8pt\kern 2pt
\refstepcounter{algorithm}
\noindent\textbf{Algorithm~\thealgorithm} FORGE generation and fail-closed synthesis-program construction.
\label{alg:forge-inference}
\par\nobreak\kern 2pt\hrule height 0.4pt\nobreak
\begin{algorithmic}[1]
\Require trained coordinate predictors, chemistry $c$, program prior $p_{\rm prog}(\cdot\mid c)$,
integers $L_{\rm samp},D\geq1$ and $B\geq0$, recorded seed $s$, inference policy $\Pi$,
route contract $\mathfrak B$
\Ensure attempt-indexed dossier sequence $\mathcal D_{\rm out}$, exact-L1 ledger $\mathcal L_{\rm out}$,
and typed failure and proposal ledger $\mathcal H$
\State initialize the random-number streams from $s$;
$\mathcal D_{\rm out}\gets[\,]$; $\mathcal L_{\rm out}\gets[\,]$; $\mathcal H\gets[\,]$
\For{$b=1,\ldots,B$}
    \State sample and record $P_b\sim p_{\rm prog}(\cdot\mid c)$; initialize raw draw ledger $\mathcal Z_b$
    \For{$d=1,\ldots,D$}
        \State sample $Z_0\sim q_0^{P_b}$ and set $Z\gets\Pi_{P_b}(Z_0)$
        \For{$k=0,\ldots,L_{\rm samp}-1$}
            \State $t_k\gets k/L_{\rm samp}$; $\Delta t_k\gets1/L_{\rm samp}$
            \State $Z\gets\Pi_{P_b}\!\left(\Call{CappedCoordinateEuler}{Z,t_k,\Delta t_k,P_b}\right)$
        \EndFor
        \State obtain terminal predictions $\lambda_{b,d}$ and $Z^{\rm out}_{b,d}\gets H_{P_b}^\theta(Z)$
        \State append the state or typed terminal failure, predictions and raw decoded graph to $\mathcal Z_b$
    \EndFor
    \State $\mathcal C_b\gets\Call{CheckedConstruction}{P_b,\mathcal Z_b,\Pi}$
    \State record every raw/proposed graph, check, abstention and cost in $\mathcal H$
\EndFor
\State $(G_1,\ldots,G_B)\gets\mathcal S_\Pi((P_b,\mathcal C_b)_{b=1}^{B})$
\For{$b=1,\ldots,B$}
    \State $G\gets G_b$; $P\gets P_b$
    \If{$G=\bot_{\rm dec}$ or $\mathrm{Valid}(G)=0$}
        \State record $(b,\textsc{InvalidGraph},\text{decode or validation reason})$ in $\mathcal H$
        \State \textbf{continue}
    \EndIf
    \State $\widehat{\mathcal{E}}_c(G)\gets
    \{(\mathbf C,\tau)\in\mathcal D_c(G):\exists H\in\mathcal F_c(\mathbf C),\ H\cong G\}$
    \If{$\widehat{\mathcal{E}}_c(G)=\varnothing$}
        \State $\mathcal H\gets\mathcal H\mathbin{\|}[(b,\textsc{L1Failure})]$
        \State \textbf{continue}
    \EndIf
    \State $\mathcal L_{\rm out}\gets\mathcal L_{\rm out}\mathbin{\|}
    [(b,G,\widehat{\mathcal E}_c(G))]$
    \If{$|\widehat{\mathcal E}_c(G)|\neq1$}
        \State $\mathcal H\gets\mathcal H\mathbin{\|}[(b,\textsc{AmbiguousL1})]$
        \State \textbf{continue}
    \EndIf
    \State select the sole $(\mathbf C_T,\tau_T)\in\widehat{\mathcal E}_c(G)$
    \State $Y\gets\mathsf{Assess}_{\mathfrak B}(G,P,\mathbf C_T,\tau_T)$
    \If{$Y$ is a dossier $T$ and $\mathrm{Complete}_{\mathfrak B}(G,T,P)=1$}
        \State $\mathcal D_{\rm out}\gets\mathcal D_{\rm out}\mathbin{\|}[(b,T)]$
    \ElsIf{$Y\in\mathfrak F_{\mathfrak B}$}
        \State $\mathcal H\gets\mathcal H\mathbin{\|}[(b,Y)]$
    \Else
        \State $\mathcal H\gets\mathcal H\mathbin{\|}[(b,\textsc{IncompleteDossier})]$
    \EndIf
\EndFor
\State \Return $\mathcal D_{\rm out},\mathcal L_{\rm out},\mathcal H$
\end{algorithmic}
\nopagebreak\hrule height 0.8pt
\endgroup

\subsection{Reproducibility and artifact provenance}
\label{app:repro}

Each run is identified by its training seed, data and split manifests, atom/program vocabularies,
source registries, fitted noise marginals, configuration, code snapshot and checkpoint SHA-256.
The independent training-seed identifiers are \PendingValue{training seed identifiers}. Each
comparison retains its declared attempt budget, checkpoint rule and evaluator. Sampling seeds,
request/layout identifiers and proposal lineage are recorded separately from training seeds.
The manuscript provenance records identify the source artifacts for tables and figures through
SHA-256 hashes.

The inference-policy identifier is \PendingValue{inference policy identifier}; its record includes
all proposal caps, TRAIN motif references, admission order, selection-group boundaries and source
checks. The route contract is \PendingValue{route contract identifier}, including depth/query
limits, evidence-admission policy and dated terminal-material snapshots. Restarts preserve committed
checkpoints and completed request shards, and source-preserving fixes retain the exact earlier
implementation and an equivalence receipt. A change to an inference map defines a distinct arm
unless equivalence is demonstrated for the declared inputs.

Reproduction requires source code, configuration and environment manifests, model checkpoints,
split/registry hashes, attempt-level output and verification ledgers, and figure/table exporters.
These artifacts are identified by \PendingValue{reproduction archive identifier}. Corpus access conditions are
\PendingValue{data availability}. Constitutional identity removes atom maps and stereochemical
labels while retaining formal charge and bond order; no tautomer normalization is applied.

\clearpage

